\appendixchapter{Lie Groups and Lie Algebras}
\label{app:lie}
\chaptermark{Lie Groups and Lie Algebras}
\section*{Introduction}
A translation, a rotation, or an invertible change of coordinates can be
composed with another transformation of the same kind. These transformations
also vary smoothly. A Lie group brings these two observations together:
\[
\boxed{\text{a Lie group is simultaneously a smooth manifold and a group}.}
\]
The compatibility of these structures transports a tangent vector at the
identity to every point. Commutators of the resulting vector fields then
encode infinitesimal noncommutativity in a finite-dimensional Lie algebra.
This is a bridge from calculus to symmetry, not a development of
representation theory or the classification of Lie groups.

We use Appendix~\ref{app:submanifolds} and only the vector-field theory
constructed below. Compare Tu~\cite[\S\S14--16]{tu-manifolds} with
Lee~\cite[Chapters 7--8,20]{lee-smooth}. The matrix examples use regular
level sets; the exponential is proved from its convergent series, avoiding
the flow-existence machinery needed for a general Lie-group exponential.
All Lie groups here are finite-dimensional real smooth manifolds without boundary.

\section{Smooth Groups and Concrete Examples}
The additive space $\R^n$ has smooth operations $(x,y)\mapsto x+y$ and
$x\mapsto-x$. The circle has complex multiplication
$(x+iy)(u+iv)=(xu-yv)+i(xv+yu)$ and inverse $x-iy$ on $x^2+y^2=1$.
These are restrictions of polynomial maps, and their images lie in the
circle. Finally invertible matrices form the open set
\[
GL_n(\R)=\{A\in M_n(\R):\det A\ne0\},
\]
because the determinant is a continuous polynomial.

\begin{definition}[Lie group]
A group is a set with an associative multiplication, an identity $e$, and
an inverse for each element. A \textbf{Lie group} is a smooth manifold $G$
with a group structure for which both maps
\[
m:G\times G\to G,\quad m(g,h)=gh,\qquad
\iota:G\to G,\quad\iota(g)=g^{-1}
\]
are smooth. A \textbf{Lie-group homomorphism} is a smooth map preserving
multiplication (and therefore identity and inverses).
\end{definition}
Thus $(\R^n,+)$ and $S^1$ are Lie groups. For $GL_n(\R)$ multiplication
has entries
\[
\sum_k a_{ik}b_{kj},
\]
 and inversion has entries given by
cofactors divided by $\det A$. They are smooth on the open set in question.
Products of Lie groups use componentwise multiplication; in particular
$T^2=S^1\times S^1$ is a Lie group with its main-text product structure.

\begin{proposition}[Special linear and orthogonal groups]\label{prop:matrix-lie-groups}
The sets $SL_n(\R)=(\det)^{-1}(\{1\})$, $O(n)=\{A:A^TA=I\}$, and
$SO(n)=\{A\in O(n):\det A=1\}$ are embedded submanifolds of matrix
space and Lie groups under matrix multiplication. Their dimensions are
$n^2-1$,
\[
n(n-1)/2,,\qquad n(n-1)/2,
\]
 respectively, for $n\geq1$.
\end{proposition}
\begin{proof}
We first verify the rank hypotheses and then restrict the operations.
Multilinearity of the determinant gives
\[
\det(I+tH)=1+t\sum_i h_{ii}+O(t^2),\qquad
D(\det)(I)H=\operatorname{tr}H,
\]
where
\[
\operatorname{tr}H=\sum_i h_{ii}.
\]
 Indeed the linear terms choose
one column from $tH$ and all other columns from $I$; only that column's
diagonal entry survives. Factoring $A+tH=A(I+tA^{-1}H)$ gives, for
invertible $A$,
\[
D(\det)(A)H=\det A\,\operatorname{tr}(A^{-1}H).
\]
At determinant $1$, choosing
\[
H=A/n
\]
 makes this derivative $1$.
The regular-level theorem therefore gives the stated dimension of $SL_n$.
The orthogonal computation, including its symmetric-matrix target and
surjectivity, is already proved in \cref{ex:orthogonal-level}.
For $A\in O(n)$, $(\det A)^2=1$, so $SO(n)$ is the open subset
$\{\det A>0\}$ of $O(n)$ and has the same local dimension.

The determinant product law and $(AB)^T(AB)=B^TA^TAB$ show that these
sets are groups; inverses remain in the same sets. The smooth operations
of $GL_n$ restrict smoothly to them by the submanifold map test in
Appendix~\ref{app:submanifolds}. This also justifies smoothness with the
induced manifold structures, rather than merely with the ambient coordinates.
\end{proof}
We have proved these examples individually. No general closed-subgroup
theorem is being assumed.

\begin{proposition}[Group operations to first order]
\label{prop:lie-first-order}
Let $m:G\times G\to G$ be multiplication and $\iota:G\to G$ inversion.
Under $T_{(e,e)}(G\times G)\cong T_eG\oplus T_eG$,
\[
dm_{(e,e)}(v,w)=v+w,\qquad d\iota_e(v)=-v.
\]
\end{proposition}
\begin{proof}
Restricting $m$ to either axis gives the identity: $m(g,e)=g=m(e,g)$.
Thus its differential sends $(v,0)$ to $v$ and $(0,w)$ to $w$;
linearity gives their sum. Differentiate $m(g,\iota(g))=e$ at $e$ to obtain
$v+d\iota_e(v)=0$.
\end{proof}
To first order, every Lie group looks like its tangent vector space with
ordinary addition. This linearization cannot detect noncommutativity.
Transporting vectors by left translation and then taking commutators of
the resulting fields supplies that extra infinitesimal information:
\[
\begin{gathered}
\text{multiplication}\longrightarrow\text{linearization}
\longrightarrow\text{left-invariant fields}\\
\longrightarrow\text{commutator}\longrightarrow\text{Lie algebra}.
\end{gathered}
\]

\section{Translation Turns One Tangent Space into Vector Fields}
\label{sec:left-invariant-fields}
For $g\in G$, left translation is $L_g(h)=gh$. It is smooth, and
$L_{g^{-1}}$ is its smooth inverse. Hence
\[
T_eG\xrightarrow{\ d(L_g)_e\ }T_gG
\]
is a linear isomorphism. One tangent space supplies directions everywhere.

A smooth vector field $X$ assigns $X_p\in T_pM$ with smooth coordinate
components. It acts on smooth scalar functions by
$Xf(p)=df_p(X_p)$. In a chart this is

\[
Xf=\sum_i X^i\partial_i f,
\]
 so $Xf$ is smooth and obeys the product rule.
\begin{definition}
A vector field $X$ on $G$ is \textbf{left invariant} if
$d(L_g)_hX_h=X_{gh}$ for all $g,h\in G$. Set $\mathfrak g=T_eG$.
\end{definition}
\begin{proposition}[Fields determined at the identity]\label{prop:left-invariant-identification}
Evaluation at $e$ is a linear bijection between smooth left-invariant
vector fields and $\mathfrak g$. Its inverse sends $v$ to the field
$X^v_g=d(L_g)_e v$.
\end{proposition}
\begin{proof}
Left invariance forces the formula, so it proves uniqueness. To check
existence including smoothness, choose a smooth curve $\gamma$ with
$\gamma(0)=e$, $\gamma'(0)=v$, using a coordinate line in a chart.
For a local target coordinate $x^j$ near $g_0$, the component of $X^v_g$
is
\[
\left.\frac{\partial}{\partial t}\right|_{t=0}x^j(g\gamma(t)).
\]
For $g$ near $g_0$ and sufficiently small $t$ the expression is defined
and smooth jointly in $(g,t)$, because multiplication is smooth. Its
$t$-derivative is smooth in $g$. Thus the field is smooth. The chain rule
and $L_g\circ L_h=L_{gh}$ give
\[
d(L_g)_hX^v_h=d(L_g)_h d(L_h)_e v=d(L_{gh})_e v=X^v_{gh}.
\]
Finally $X^v_e=v$, and linearity follows from that of each differential.
\end{proof}
For a matrix group the transported field is $X^A_g=gA$. For the additive
group it is the constant field $X^v_x=v$. These formulas are concrete
realizations of the same transport mechanism.

\section{The Bracket and the Lie Algebra}
Composition of vector fields acting on functions generally has second
derivatives. Subtracting the reverse composition cancels them, leaving a
first-order operation and hence another vector field.
\begin{proposition}[Vector-field commutator]\label{prop:vector-field-bracket}
For smooth vector fields $X,Y$, the rule
\[
[X,Y]f=X(Yf)-Y(Xf)
\]
defines a smooth vector field, with chart components
\[
[X,Y]^j=\sum_i\bigl(X^i\partial_iY^j-Y^i\partial_iX^j\bigr).
\]
The bracket is bilinear, antisymmetric, and satisfies the Jacobi identity
$[X,[Y,Z]]+[Y,[Z,X]]+[Z,[X,Y]]=0$.
\end{proposition}
\begin{proof}
Expanding the first composition gives
\[
X(Yf)=\sum_{i,j}X^i(\partial_iY^j)\partial_jf
+\sum_{i,j}X^iY^j\partial_i\partial_jf.
\]
Interchange $X,Y$ in the second composition. Reindex $i,j$ and use equality
of mixed partials to cancel the second-derivative terms. The remaining
coefficient is the stated smooth vector field. This construction is
intrinsic because its action on every local scalar function is intrinsic;
coordinate functions determine each tangent vector. Bilinearity and
antisymmetry follow from the subtraction. For Jacobi, the three expansions
as operators are
\[
\begin{aligned}
[X,[Y,Z]]&=XYZ-XZY-YZX+ZYX,\\
[Y,[Z,X]]&=YZX-YXZ-ZXY+XZY,\\
[Z,[X,Y]]&=ZXY-ZYX-XYZ+YXZ.
\end{aligned}
\]
Every term cancels with its opposite; evaluation on local coordinates
then gives the zero vector field.
\end{proof}
If $h:M\to N$ is a diffeomorphism, define
$(h_*X)_{h(p)}=dh_pX_p$. The chain rule gives
$((h_*X)f)\circ h=X(f\circ h)$. Applying this identity twice and
subtracting proves
\[
h_*[X,Y]=[h_*X,h_*Y].
\]
In particular the bracket of left-invariant fields is left invariant:
take $h=L_g$ and use the defining invariances.

\begin{definition}[Lie algebra of a Lie group]
A real \textbf{Lie algebra} is a real vector space with a bilinear
antisymmetric bracket satisfying the Jacobi identity. On $\mathfrak g=T_eG$
define
\[
[v,w]_{\mathfrak g}=[X^v,X^w]_e.
\]
\end{definition}
By \cref{prop:left-invariant-identification}, the left-invariant extension
of this bracket is $[X^v,X^w]$. Thus all three identities pass from vector
fields to $\mathfrak g$, proving that this definition gives a Lie algebra.
Its underlying vector space has dimension $\dim G$.

\begin{proposition}[Matrix Lie algebras]\label{prop:matrix-lie-algebras}
With tangent spaces identified in matrix space,
\[
\begin{aligned}
\mathfrak{gl}_n&=M_n(\R),\\
\mathfrak{sl}_n&=\{A:\operatorname{tr}A=0\},\\
\mathfrak{so}(n)&=\{A:A^T=-A\}.
\end{aligned}
\]
Both $O(n)$ and $SO(n)$ have the last tangent space at $I$.
For these matrix groups the bracket is $[A,B]=AB-BA$.
\end{proposition}
\begin{proof}
The first group is open. For $SL_n$, use the kernel of
$D(\det)(I)H=\operatorname{tr}H$. For $O(n)$ use the kernel of
$H\mapsto H+H^T$; the open subset $SO(n)$ has the same tangent space.
For the bracket, the left-invariant fields are $X^A(g)=gA$ and
$X^B(g)=gB$. The derivatives of these ambient linear expressions are
$H\mapsto HA$ and $H\mapsto HB$. The coordinate bracket formula gives
\[
[X^A,X^B](g)=(gA)B-(gB)A=g(AB-BA).
\]
The calculation is also valid after restriction to an embedded matrix
group: on every ambient matrix coordinate the restricted fields act by
the same derivative, and these coordinates determine the included tangent
vector. Evaluation at $I$ proves the claim. Closure can also be checked
directly: $\operatorname{tr}(AB)=\operatorname{tr}(BA)$ by interchanging
the two summation indices, and for skew $A,B$,
$(AB-BA)^T=BA-AB$.
\end{proof}

\begin{example}[A nonzero bracket]\label{ex:nonzero-lie-bracket}
In $\mathfrak{sl}_2$, take
\[
H=\begin{pmatrix}1&0\\0&-1\end{pmatrix},\qquad
E=\begin{pmatrix}0&1\\0&0\end{pmatrix}.
\]
Then $HE=E$ and $EH=-E$, so the bracket from
\cref{prop:matrix-lie-algebras} is $[H,E]=2E\ne0$.
Equivalently the left-invariant fields satisfy
$[X^H,X^E](g)=2gE$: their commutator, as defined in
\cref{prop:vector-field-bracket}, need not vanish.
\end{example}

\begin{remark}[The principal examples together]
The tangent descriptions above give the following comparison. The circle
and torus use their usual angle coordinates at the identity.
\[
\begin{array}{c|c|c}
G&T_eG&\text{bracket}\\\hline
(\R^n,+)&\R^n&0\\
S^1&\R&0\\
T^n&\R^n&0\\
GL_n(\R)&M_n(\R)&AB-BA\\
SL_n(\R)&\{A:\operatorname{tr}A=0\}&AB-BA\\
O(n)&\{A:A^T=-A\}&AB-BA\\
SO(n)&\{A:A^T=-A\}&AB-BA
\end{array}
\]
For the additive group the invariant fields are constant, so their bracket
vanishes. Angle coordinates give the same local calculation for circles
and tori. Equal Lie algebras need not mean equal global groups.
\end{remark}

\section{Matrix Exponentials and One-Parameter Subgroups}
\label{sec:matrix-exponential}
A tangent matrix describes initial velocity. Its exponential integrates
that velocity into a family of transformations. We can prove this directly
for matrices with the convergence tools already available.
\begin{definition}
For $A\in M_n(\R)$ set
\[
\exp(A)=\sum_{k=0}^{\infty}\frac{A^k}{k!},\qquad A^0=I.
\]
The notation $e^A$ is a standard alternative for this series-defined
matrix exponential; it is not an application of real-power laws to matrices.
A \textbf{one-parameter subgroup} of a Lie group $G$ is a smooth map
$\gamma:\R\to G$ with $\gamma(t+s)=\gamma(t)\gamma(s)$.
\end{definition}
\begin{proposition}[The exponential laws needed here]
The series converges, and $\gamma(t)=\exp(tA)$ is smooth with
\[
\begin{aligned}
\gamma(0)&=I,& \gamma'(t)&=A\exp(tA)=\exp(tA)A,\\
\exp(tA)\exp(sA)&=\exp((t+s)A),&
\exp(tA)^{-1}&=\exp(-tA).
\end{aligned}
\]
Every one-parameter subgroup of $GL_n(\R)$ has this form, with
$A=\gamma'(0)$.
\end{proposition}
\begin{proof}
Use an operator norm, so $\|A^k\|\leq\|A\|^k$. The scalar exponential
series gives absolute convergence and uniform convergence for bounded $t$.
For each derivative order the differentiated series is bounded on compact
$t$-intervals by a constant times an exponential series. The termwise
differentiation theorem therefore gives smoothness and the derivative
formula. For the product, absolute convergence permits the Cauchy product:
\[
\begin{aligned}
\exp(tA)\exp(sA)
&=\sum_{m=0}^{\infty}\left(\sum_{k=0}^m
\frac{t^ks^{m-k}}{k!(m-k)!}\right)A^m\\
&=\sum_{m=0}^{\infty}\frac{(t+s)^mA^m}{m!}.
\end{aligned}
\]
Taking $s=-t$ proves invertibility. Conversely differentiate
$\gamma(t+s)=\gamma(t)\gamma(s)$ in $s$ at zero to get
$\gamma'(t)=\gamma(t)A$. The product rule gives
$(\gamma(t)\exp(-tA))'=0$, hence $\gamma(t)\exp(-tA)=I$ by its value at zero.
This proves the converse without an ODE existence theorem.
\end{proof}
The same exponential law is not asserted for $\exp(A)\exp(B)$ with arbitrary
noncommuting $A,B$.

\begin{proposition}[Exponentiating the displayed Lie algebras]
If $\operatorname{tr}A=0$, then $\exp(tA)\in SL_n(\R)$. If $A^T=-A$,
then $\exp(tA)\in SO(n)$. Thus the exponential supplies one-parameter
subgroups with every initial tangent vector for the matrix groups above.
\end{proposition}
\begin{proof}
Set $u(t)=\det(\exp(tA))$. The determinant derivative and the preceding
proposition give $u'(t)=u(t)\operatorname{tr}A$, $u(0)=1$.
Differentiating $e^{-t\operatorname{tr}A}u(t)$ proves
$u(t)=e^{t\operatorname{tr}A}$. For skew $A$, transposing the series gives
$(\exp(tA))^T=\exp(-tA)$, so the matrix is orthogonal. Its determinant is $1$,
either by the trace formula or by continuity from $t=0$ into $\{1,-1\}$.
\end{proof}
\begin{example}[The whole computation for $SO(2)$]
Every skew $2$ by $2$ matrix is $A=\theta J$, where
\[
J=\begin{pmatrix}0&-1\\1&0\end{pmatrix},\qquad J^2=-I.
\]
Separate the even and odd powers in the absolutely convergent series:
\[
\exp(\theta J)=\cos\theta\,I+\sin\theta\,J
=\begin{pmatrix}\cos\theta&-\sin\theta\\
\sin\theta&\cos\theta\end{pmatrix}.
\]
The one-parameter subgroup $t\mapsto \exp(t\theta J)$ rotates at constant
angular speed $\theta$. Acting on $(1,0)$ traces the unit circle;
the initial velocity is $\theta J(1,0)=(0,\theta)$, its tangent vector
at that point. More generally $J(\cos t,\sin t)=(-\sin t,\cos t)$
is the counterclockwise unit tangent. Thus the matrix generator, circle
tangent, and exponential describe the same motion. Here $\mathfrak{so}(2)=\R J$ has zero bracket,
and exponentiation repeats after $2\pi J$. Infinitesimal data do not by
themselves distinguish the additive line from the circle's periodic global
geometry. For general Lie groups an exponential is constructed by integral
curves of left-invariant fields; that existence theory is beyond this appendix.
\end{example}

\par\Needspace{12\baselineskip}
\section*{Exercises}
\addcontentsline{toc}{section}{Exercises}
\begin{hexercise}{app-g-1}
Check the group and smoothness axioms for $S^1$ using real coordinates.
Identify $T_{(1,0)}S^1$ and its left-invariant fields.
\end{hexercise}
\begin{hexercise}{app-g-2}
Compute $T_ASL_n$ and $T_AO(n)$ at a general point and show they are
obtained from their identity tangent spaces by $H\mapsto AH$.
\end{hexercise}
\begin{hexercise}{app-g-3}
For $X=x\partial_y$ and $Y=y\partial_x$ on $\R^2$, compute $[X,Y]$.
Verify its action directly on $x$, $y$, and $xy$.
\end{hexercise}
\begin{hexercise}{app-g-4}
In $\mathfrak{sl}_2$, put $H=\operatorname{diag}(1,-1)$,
$E=\left(\begin{smallmatrix}0&1\\0&0\end{smallmatrix}\right)$, and
$F=\left(\begin{smallmatrix}0&0\\1&0\end{smallmatrix}\right)$.
Compute all pairwise brackets and verify one Jacobi identity directly.
\end{hexercise}
\begin{hexercise}{app-g-5}
Compute $\exp(tA)$ for a diagonal matrix and for
$A=\left(\begin{smallmatrix}0&1\\0&0\end{smallmatrix}\right)$.
Check the subgroup law and the initial velocity in each case.
\end{hexercise}
\begin{hexercise}{app-g-6}
Show that a smooth homomorphism $\gamma:\R\to S^1$ has the form
$\gamma(t)=(\cos at,\sin at)$ by embedding the circle as $SO(2)$ and
using the one-parameter subgroup proposition.
\end{hexercise}
\begin{exercise}
Show that $O(n)$ and $SO(n)$ have identical Lie algebras although they
are different groups for $n\geq1$. Identify explicitly which global
information their tangent spaces at the identity fail to record.
\end{exercise}
